\documentclass[border=8pt]{standalone}
\usepackage{polymer-paper-preamble}

%% Smoke test: 2-layer variant of isometric_layered_device_stack
%% Demonstrates composition freedom: remove Layer B (dielectric) and Layer D (top electrode)
%% Keep: Layer A (bottom) + Layer C (charge)

\begin{document}
\resizebox{92mm}{!}{%
\begin{tikzpicture}[iso basis]
  \node[font=\fontsize{9}{10.8}\selectfont\bfseries, text=cGray!90!black, anchor=west]
    at (-0.40, -1.10, 1.95) {Simplified: substrate + charge};

  % Layer A: substrate
  \begin{scope}[shift={(0,0,0)}]
    \IsoBlock{4.8}{1.7}{0.28}{cAmber}{}
    \node[iso label, anchor=east] at (-0.20,0.10,0.16) {substrate};
    \draw[iso callout] (-0.10,0.12,0.16) -- (0.65,0.10,0.16);
  \end{scope}

  % Layer C: charge region (moved up, no intermediate layers)
  \begin{scope}[shift={(0,0,0.40)}]
    \IsoBlock{4.8}{1.7}{0.34}{cBlue}{}
    \foreach \x/\y in {0.80/0.45,1.50/1.10,2.25/0.72,3.05/1.28}{
      \fill[iso charge] (\x,\y,0.20) circle[radius=0.055];
    }
    \node[iso label, anchor=west] at (5.95,-0.70,0.28) {interface charges};
    \draw[iso callout] (5.75,-0.64,0.28) -- (3.78,0.58,0.20);
  \end{scope}

  % Depth cue
  \draw[iso hidden edge] (4.95,1.90,0.28) -- (4.95,1.90,0.92);
  \node[iso label, text=cGray!70, anchor=west] at (5.95,-0.50,0.60) {fixed-view depth cue};
\end{tikzpicture}%
}
\end{document}
